% Generated draft: requires mathematical and source review.
\documentclass[11pt]{article}
\usepackage[margin=25mm]{geometry}
\usepackage{fontspec}
\setmainfont{Latin Modern Roman}
\usepackage{amsmath,amssymb,mathtools,unicode-math}
\setmathfont{Latin Modern Math}
\usepackage{enumitem,needspace,xurl,hyperref}
\setlength{\parindent}{0pt}
\setlength{\parskip}{5pt}
\setlength{\emergencystretch}{3em}
\newcommand{\N}{\mathbb{N}}
\newcommand{\Z}{\mathbb{Z}}
\newcommand{\Q}{\mathbb{Q}}
\newcommand{\R}{\mathbb{R}}
\newcommand{\C}{\mathbb{C}}
\newcommand{\F}{\mathbb{F}}
\newcommand{\A}{\mathbb{A}}
\newcommand{\PP}{\mathbb{P}}
\newcommand{\E}{\mathbb{E}}
\newcommand{\eps}{\varepsilon}
\newcommand{\dd}{\,\mathrm{d}}
\begin{document}
\section*{UnsolvedMath ID: 30001553}
\textbf{Source problem code:} catalogue entry
\subsection{Definitions and mathematical statement}
Let \(A\) be an alphabet and let \(A^*\) denote the set of finite words over \(A\), including the empty word \(\varepsilon\). The source does not specify whether \(A\) is finite. A morphic involution is a map \(\theta:A^*\to A^*\) satisfying
\[
\theta(\varepsilon)=\varepsilon,\qquad
\theta(uv)=\theta(u)\theta(v),\qquad
\theta(\theta(u))=u
\quad (u,v\in A^*).
\]
For a word \(w\), its length is denoted by \(|w|\), and a contiguous factor is a word \(u\) for which \(w=xuy\) for some \(x,y\in A^*\).

The supplied record describes \(\tau_\theta(w)\) as the maximum length of a \(\theta\)-unbordered factor of \(w\), and \(\pi_\theta^{\mathrm{alt}}(w)\) as the shortest alternating \(\theta\)-period of \(w\). However, it supplies neither a definition of a \(\theta\)-border nor a definition of an alternating \(\theta\)-period. Consequently these descriptions do not determine the two quantities precisely. In particular, the admissible factors and period lengths, the treatment of partial terminal blocks, and the conventions for empty words or empty sets of candidates remain unspecified.

\textbf{Source problem, with unresolved terminology retained.} Determine the sharp relation between \(\tau_\theta(w)\) and \(\pi_\theta^{\mathrm{alt}}(w)\) for a morphic involution \(\theta\) on \(A^*\), analogous to the ordinary relation between the maximum length of an unbordered factor and the shortest period. The source does not identify the ordinary relation intended as the analogy, state a proposed inequality or equality, or specify whether additional parameters such as \(|w|\) are permitted. The intended universal scope over alphabets, involutions, and words also requires clarification. This is therefore a statement of the supplied research question with explicit gaps, rather than a fully specified conjecture.
\subsection{Short English statement}
The problem asks for a sharp comparison between the longest \(\theta\)-unbordered factor and the shortest alternating \(\theta\)-period of a finite word. The supplied record omits the definitions needed to make that comparison precise, including boundary conventions and the intended ordinary-period analogy.
\subsection{Sources}
\begin{itemize}
\item Local UnsolvedMath registry: supplied complete record, permanent ID \(30001553\), source problem code \texttt{catalogue entry}, title \emph{Sharp Bounds for Alternating Involution Periods}, in the Oberwolfach Reports Open Problems collection. No individual proposer is supplied; the proposed year is recorded as 2010.
\item Supplied source citation: \emph{Mini-Workshop: Combinatorics on Words (2011)}, Oberwolfach Reports, pp.~2195--2244, DOI \texttt{10.4171/owr/2010/37}. Original workshop report: \url{https://doi.org/10.4171/owr/2010/37}. The record describes Oberwolfach Reports as documenting mathematical workshops and their open-problem discussions. The citation's 2011 label and the record's 2010 proposed year are retained without reconciliation.
\item Supplied literature source: Kari and Kulkarni, \url{https://doi.org/10.1007/s00236-016-0258-1}. No title or specific theorem supporting this formulation is supplied.
\item Supplied literature assessment, checked on \textbf{2026-08-21}: the record labels the problem open and reports that later work establishes special properties of pseudo-bordered words, while its search found no general resolution of the relation between unbordered and \(\theta\)-bordered factors. This is supplied dated curation derived from the OWR material, not independently verified here; its status and cited claims require independent verification.
\item The record labels the statement \texttt{corrected\_verified}, but its accompanying verification note reports only enclosing TeX notation in mathematical delimiters without changing mathematical content. This does not establish the missing definitions or mathematical correctness.
\end{itemize}
\end{document}
